% =========================
% Cas pratique 9
% =========================

\section[Cas 9]{Cas pratique 9 : Droit de possession
personnelle antérieure}

%--------------------------
\begin{frame}{Question 1}

\begin{block}{}
L'article \textbf{L613-3}(a) CPI (actes de contrefaçon)
+ \textbf{L615-1 al. 3} CPI (conditions de responsabilité)
+ \textbf{L615-4} CPI (opposable à partir de la publication).
\end{block}

\begin{itemize}
  \item En l'espèce, Padrole a un brevet \textbf{français} delivré.
  \begin{itemize}
    \item L'élément matériel: la revendication de produit sur la
  pince.
    \item L'élément légal: acte de fabrication de Mercurey,
  acte d'offre, de mise sur le marché et détention en vue
  de l'offre et de la mise sur le marché de Langsamer.
  Padrole bénéficie de la connaissance de cause, mais
  pas Mercurey.
    \item L'élément temporel: publication de la demande de brevet
  en août 2015. Les actes de contrefaçon ne sont comptabilisés
  qu'à partir de ce mois-là. Or, la commercialisation a
  commencé en janvier 2015. 
  \end{itemize}
  \item En conclusion, un risque de contrefaçon pour Padrole
  et Mercurey existe sur le territoire français.
\end{itemize}

\end{frame}

%--------------------------
\begin{frame}{Question 2}

\begin{block}{}
L'article \textbf{L611-11 CPI} (nouveauté)
+ \textbf{L611-13 CPI} (état de la technique).
\end{block}

\begin{itemize}
  \item En l'espèce, les prototypes dès décembre 2013 < dépôt FR.
  Mais une clause de confidentialité.
  \item En conclusion, raisonnement sur défaut de nouveauté
  pas de fondement. A moins d'une diffusion illicite
  prouvée avant la date de dépôt.
\end{itemize}

\end{frame}

%--------------------------
\begin{frame}{Question 3}

\begin{block}{}
L'article \textbf{L613-7 CPI} (droit
de possession personnelle antérieure) dispose
d'une exception à L613-3 CPI.
\begin{itemize}
  \item une possession de l’invention sur le territoire français,
  \item une exploitation effective ou des préparatifs sérieux,
  \item une possession de bonne foi.
\end{itemize}
\end{block}  

\begin{itemize}
  \item En l'espèce, \begin{itemize}
    \item Mercurey possède l'invention, pas Langsamer.
    \item 3 ans de R\&D.
    \item Pas d'élément de mauvaise foi.
  \end{itemize}
  \item En conclusion, Mercurey, pas Langsamer, a le droit de se
  prévaloir d'une possession personnelle antérieure.
\end{itemize}

\end{frame}